% Part: incompleteness
% Chapter: arithmetization-syntax
% Section: proofs-in-nd

\documentclass[../../../include/open-logic-section]{subfiles}

\begin{document}

\olfileid{inc}{art}{pnd}
\olsection{فطری استنباط وچ \usetoken{P}{derivation}}

\begin{explain}
!!{derivation}s نوں عددی روپ دین لئی سانوں !!{derivation}s نوں
عدداں نال ظاہر کرنا پینا اے۔ کیوں جے !!{derivation}s،
!!{formula}s دے درخت نیں تے ہر استنباط نال اک یا دو لیبل وی
ہوندے نیں، ایس لئی بازگشتی طریقہ سِدھا اے۔
اسیں !!{derivation} نوں اک ٹپل نال ظاہر کردے آں۔
اوس ٹپل دے جز آخری استنباط دے مقدمیاں تک پہنچن والیاں
براہِ راست ذیلی !!{derivation}s دی گنتی، اوہناں ذیلی !!{derivation}s دے عددی روپ، تے آخرلا !!{formula}،
آخری استنباط دا مفروضہ خارج کرن والا لیبل، تے آخری
استنباط دی قسم دسن والا اک عدد نیں۔
\end{explain}

\begin{defn}
  جے $\delta$ فطری استنباط وچ !!a{derivation} اے،
  تاں $\Gn{\delta}$ دی استقرائی تعریف ایویں کردے آں:
  \begin{enumerate}
  \item
    جے $\delta$ صرف مفروضے~$!A$ اُتے مشتمل اے،
    تاں $\Gn{\delta}$، $\tuple{0, \Gn{!A}, n}$ اے۔
    جے ایہ !!{undischarged} مفروضہ اے تاں عدد~$n$،
    $0$ اے؛ ورنہ اوس دا عددی لیبل اے۔
  \item
    جے $\delta$ دا آخرلا استنباط صفر، اک، دو یا
    تن مقدمیاں والا اے، تاں $\Gn{\delta}$ بالترتیب
    ایہ اے:
    \begin{align*}
      & \tuple{0, \Gn{!A}, n, k}, \\
      & \tuple{1, \Gn{\delta_1}, \Gn{!A}, n, k}, \\
      & \tuple{2, \Gn{\delta_1}, \Gn{\delta_2}, \Gn{!A}, n, k}, \text{ یا}\\
      & \tuple{3, \Gn{\delta_1}, \Gn{\delta_2}, \Gn{\delta_3}, \Gn{!A}, n,
        k},
    \end{align*}
    ایتھے $\delta_1$، $\delta_2$ تے $\delta_3$،
    $\delta$ دے آخری استنباط دے مقدمے یا مقدمیاں
    اُتے ختم ہون والے ذیلی !!{derivation}s نیں۔
    $!A$، $\delta$ وچ آخری استنباط دا نتیجہ اے؛
    $n$ اوس استنباط دا مفروضہ خارج کرن والا لیبل
    اے (جے استنباط کوئی مفروضہ خارج نہیں کردا
    تاں $0$)؛ تے آخری استنباط وچ ورتے قاعدے مطابق $k$ دی قدر ہِٹھلے جدول توں ملدی اے۔

    \begin{tabular}{lccccccc}
      \text{قاعدہ:} & \Intro{\land} & \Elim{\land} & \Intro{\lor} & \Elim{\lor} \\
      $k$: & 1 & 2 & 3 & 4  \\[1.5ex]
      \text{قاعدہ:} &  \Intro{\lif} & \Elim{\lif} & \Intro{\lnot} & \Elim{\lnot} \\
      $k$: & 5 & 6 & 7 & 8 \\[1.5ex]
      \text{قاعدہ:}  &  \FalseInt & \FalseCl & \Intro{\lforall} & \Elim{\lforall} \\
      $k$: & 9 & 10 & 11 & 12 \\[1.5ex]
      \text{قاعدہ:} & \Intro{\lexists} & \Elim{\lexists} & \Intro{\eq} & \Elim{\eq} \\
      $k$: & 13 & 14 & 15 & 16
    \end{tabular}
  \end{enumerate}
\end{defn}

\begin{ex}
  ایہ سوکھا جیہا !!{derivation} ویکھو:
  \begin{prooftree}
    \AxiomC{$\Discharge{!A \land !B}{1}$}
    \RightLabel{\Elim{\land}}
    \UnaryInfC{$!A$}
    \DischargeRule{\Intro{\lif}}{1}
    \UnaryInfC{$(!A \land !B) \lif !A$}
  \end{prooftree}
  مفروضے دا گوڈل نمبر
  $d_0 = \tuple{0, \Gn{!A \land !B}, 1}$ ہوئے گا۔
  $\Elim{\land}$ دے نتیجے اُتے ختم ہون والے
  !!{derivation} دا گوڈل نمبر
  $d_1 = \tuple{1, d_0, \Gn{!A}, 0, 2}$
  ہوئے گا ($1$ ایس لئی کہ $\Elim{\land}$ دا اک مقدمہ
  اے؛ نتیجے~$!A$ دا گوڈل نمبر؛ $0$ ایس لئی کہ
  کوئی مفروضہ خارج نہیں ہوندا؛ تے $2$،
  $\Elim{\land}$ نوں کوڈ کرن والا عدد اے)۔
  فیر پورے !!{derivation} دا گوڈل نمبر
  $\tuple{1, d_1, \Gn{((!A \land !B) \lif
      !A)}, 1, 5}$ اے؛ یعنی
  \[
  \tuple{1, \tuple{1, \tuple{0, \Gn{(!A \land !B)}, 1}, \Gn{!A}, 0, 2},
    \Gn{((!A \land !B) \lif !A)}, 1, 5}.
  \]
\end{ex}

\begin{explain}
!!{derivation}s دی عددی روپ طے کر کے
سانوں ایہ وی دکھانا پینا اے کہ اسیں انہاں
!!{derivation}s دے گوڈل نمبراں اُتے بنیادی
بازگشتی ڈھنگ نال عمل کر سکدے آں، تے اوہناں دیاں ضروری خاصیتاں تے تعلقاں نوں وی ایس طریقے نال بیان کر سکدے آں۔ کجھ عمل سادے نیں: مثال لئی،
!!a{derivation} دے گوڈل نمبر~$d$ توں
$\fn{EndFmla}(d) = (d)_{(d)_0+1}$
اوس دے آخرلے !!{formula} دا گوڈل نمبر
دیندا اے؛
$\fn{DischargeLabel}(d) = (d)_{(d)_0 + 2}$
مفروضہ خارج کرن والا لیبل دیندا اے؛ تے
$\fn{LastRule}(d) = (d)_{(d)_0 + 3}$
آخری استنباط دی قسم دسن والا عدد دیندا اے۔
کجھ ہور عمل کافی اوکھے نیں۔ اسیں گھٹ توں
گھٹ انہاں دا طریقہ خاکہ وار دساں گے۔
مقصد ایہ دکھانا اے کہ
«$\delta$،~$!A$ دا
!!a{derivation} $\Gamma$ توں اے» والا تعلق،
$\delta$ تے~$!A$ دے گوڈل نمبراں
دا بنیادی بازگشتی تعلق اے۔
\end{explain}

\begin{prop}
  ہِٹھلے تعلق بنیادی بازگشتی نیں:
  \begin{enumerate}
  \item $!A$، $\delta$ وچ لیبل~$n$ والا
    مفروضے وجوں آؤندا اے۔
  \item $\delta$ وچ لیبل~$n$ والے سارے
    مفروضے $!A$ دے روپ وچ نیں
    (یعنی اسیں مفروضہ~$!A$ نوں
    لیبل~$n$ ورت کے $\delta$ وچ
    !!{discharge} کر سکدے آں)۔
  \end{enumerate}
\end{prop}

\begin{proof}
  سانوں دکھانا اے کہ !!{formula}s دے گوڈل
  نمبراں تے !!{derivation}s دے گوڈل نمبراں
  دے وچکار متعلقہ تعلق بنیادی بازگشتی نیں۔
  \begin{enumerate}
  \item اسیں دکھانا چاہندے آں کہ
    $\fn{Assum}(x, d, n)$ بنیادی بازگشتی اے۔
    ایہ اوس ویلے سچ ہوندا اے جدوں $x$،
    گوڈل نمبر~$d$ والے !!{derivation}
    دے لیبل~$n$ والے کسے مفروضے دا گوڈل
    نمبر ہووے۔ ایہ اوس ویلے اے جدوں
    گوڈل نمبر~$\tuple{0, x, n}$ والا
    !!{derivation}، $d$ دا ذیلی
    !!{derivation} ہووے۔ دھیان دیو کہ
    !!{derivation}s دی ساڈی کوڈکاری،
    \olref[cmp][rec][tre]{sec} وچ دتی
    درختاں دی کوڈکاری دی اک خاص صورت اے۔
    ایس لئی بنیادی بازگشتی فنکشن
    $\fn{SubtreeSeq}(d)$، $d$ دے سارے
    ذیلی !!{derivation}s دے گوڈل نمبراں
    دا سلسلہ دیندا اے (جیہدی لمبائی ودھ توں ودھ $d$ اے)۔
    سو اسیں ایہ تعریف کر سکدے آں:
    \[
    \fn{Assum}(x, d, n) \defiff \bexists{i<d}{(\fn{SubtreeSeq}(d))_i =
      \tuple{0, x, n}}.
    \]
  \item اسیں دکھانا چاہندے آں کہ
    $\fn{Discharge}(x, d, n)$ بنیادی
    بازگشتی اے۔ ایہ اوس ویلے سچ ہوندا
    اے جدوں لیبل~$n$ والے،
    گوڈل نمبر~$d$ دے
    !!{derivation} وچلے سارے
    مفروضے، گوڈل نمبر~$x$
    والے !!{formula} ہی ہون۔
    ایہ تعلق عین اوس ویلے سچ اے جدوں
    $\bforall{y<d}{(\fn{Assum}(y, d, n) \lif y
      = x)}$ ہووے۔
  \end{enumerate}
\end{proof}

\begin{prop}
  \ollabel{prop:followsby}
  خاصیت $\fn{Correct}(d)$ اُدوں تے صرف اُدوں قائم اے جدوں
!!{derivation}~$\delta$ دا آخری استنباط درست ہووے، جتھے اوس اخذ
دا گوڈل نمبر~$d$ اے۔ ایہ خاصیت بنیادی بازگشتی اے۔
\end{prop}

\begin{proof}
  ایتھے سانوں استنباط دے ہر قاعدے~$R$
  لئی ایہ دکھانا پینا اے کہ تعلق
  $\fn{FollowsBy}_R(d)$ بنیادی بازگشتی
  اے۔ $\fn{FollowsBy}_R(d)$ عین اوس
  ویلے سچ ہوندا اے جدوں $d$ کسے
  !!{derivation}~$\delta$ دا گوڈل نمبر
  ہووے، تے $\delta$ دا آخرلا
  !!{formula}، قاعدہ~$R$ ٹھیک طرح لاؤن نال $\delta$ دے براہِ راست ذیلی
!!{derivation}s دے نتیجیاں توں نکلدا ہووے۔

  اک سوکھی مثال \Intro{\land} دا قاعدہ اے۔ جے $\delta$ اک درست
$\Intro{\land}$ استنباط اُتے ختم ہووے، تاں اوہ ایویں دِسے گا:
  \begin{prooftree}
    \AxiomC{}
    \RightLabel{$\delta_1$}
    \DeduceC{$!A$}

    \AxiomC{}
    \RightLabel{$\delta_2$}
    \DeduceC{$!B$}

    \RightLabel{\Intro\land}
    \BinaryInfC{$!A \land !B$}
  \end{prooftree}
  فیر گوڈل نمبر~$d$، جیہڑا~$\delta$ نوں دِتا گیا اے،
  $\tuple{2, d_1, d_2,
    \Gn{(!A \land !B)}, 0, k}$ اے؛ جتھے
  $\fn{EndFmla}(d_1) = \Gn{!A}$،
  $\fn{EndFmla}(d_2) = \Gn{!B}$،
  $n=0$ تے $k=1$ نیں۔ سو اسیں
  $\fn{FollowsBy}_{\Intro\land}(d)$
  نوں ایس طرح تعریف کر سکدے آں:
  \begin{multline*}
    (d)_0 = 2 \land \fn{DischargeLabel}(d) = 0 \land \fn{LastRule}(d) = 1 \land {}\\
    \fn{EndFmla}(d) = {}\\
    \Gn{(} \concat \fn{EndFmla}((d)_1) \concat \Gn{\land}
    \concat \fn{EndFmla}((d)_2) \concat \Gn{)}.
  \end{multline*}

  اک ہور سادہ مثال $\Intro\eq$ دا قاعدہ
  اے۔ ایس دے کوئی مقدمے نہیں، ایس لئی
  $(d)_0 = 0$ اے، مفروضیاں وانگ۔
  ایس دا مفروضہ خارج کرن والا لیبل وی
  نہیں؛ یعنی $n=0$۔ پر $!A$ ضرور
  $\eq[t][t]$ دے روپ وچ ہووے،
  جتھے~$t$ بند ٹرم اے۔
  ایتھے اک بنیادی بازگشتی تعریف ایہ اے:
  \begin{multline*}
    (d)_0 = 0 \land \fn{DischargeLabel}(d) = 0 \land {}\\
    \bexists{t<d}{(\fn{ClTerm}(t) \land
      \fn{EndFmla}(d) = {}\\
      \Gn{{\eq}(} \concat t \concat \Gn{,} \concat t \concat \Gn{)})}.
  \end{multline*}

  ذرا اوکھی مثال وچ
  $\fn{FollowsBy}_{\Intro{\lif}}(d)$
  عین اوس ویلے سچ ہوندا اے جدوں
  $\delta$ دا آخرلا !!{formula}،
  $(!A \lif !B)$ دے روپ وچ ہووے،
  $\delta_1$ دا آخرلا !!{formula}
  $!B$ ہووے، تے $\delta$ وچ
  لیبل~$n$ والا ہر مفروضہ
  $!A$ دے روپ وچ ہووے۔
  اسیں ایس نوں بنیادی بازگشتی
  طریقے نال ایویں بیان کر
  سکدے آں:
  \begin{multline*}
    (d)_0 = 1 \land {}\\
    \bexists{a<d}{(\fn{Discharge}(a, (d)_1, \fn{DischargeLabel}(d)) \land {}}\\
      \fn{EndFmla}(d) = (\Gn{(} \concat a \concat \Gn{\lif}
      \concat \fn{EndFmla}((d)_1) \concat \Gn{)}))
  \end{multline*}
  ($a$ نوں~$!A$ دا گوڈل نمبر سمجھو)۔

  اک ہور مثال لئی \Intro{\lexists}
  نوں ویکھو۔ ایتھے $\delta$
  دا آخری استنباط عین اوس
  ویلے درست اے جدوں کوئی
  !!a{formula}~$!A$، بند ٹرم~$t$
  تے !!a{variable}~$x$ ایہوجیہے
  ہون کہ $\Subst{!A}{t}{x}$،
  !!{derivation}~$\delta_1$ دا
  آخرلا !!{formula} ہووے تے
  $\lexists[x][!A]$ آخری
  استنباط دا نتیجہ ہووے۔
  سو $\fn{FollowsBy}_{\Intro{\lexists}}(d)$
  عین اوس ویلے سچ ہوندا اے جدوں
  \begin{multline*}
    (d)_0 = 1 \land \fn{DischargeLabel}(d) = 0 \land {} \\
    \bexists{a < d}{\bexists{x<d}{\bexists{t<d}{
          (\fn{ClTerm}(t) \land \fn{Var}(x)  \land {}}}}\\
    \fn{Subst}(a,t,x) = \fn{EndFmla}((d)_1) \land
    \fn{EndFmla}(d) = (\Gn{\lexists} \concat x \concat a)).
  \end{multline*}

  فیر اسیں $\fn{Correct}(d)$ دی تعریف ایویں کردے آں:
  \begin{multline*}
    \fn{Sent}(\fn{EndFmla}(d)) \land {}\\
    (\fn{LastRule}(d) = 1 \land
    \fn{FollowsBy}_{\Intro\land}(d)) \lor \dots \lor {}\\
    (\fn{LastRule}(d) = 16 \land \fn{FollowsBy}_{\Elim\eq}(d)) \lor {}\\
    \bexists{n<d}{\bexists{x<d}{(d = \tuple{0, x, n})}}.
  \end{multline*}
  پہلی سطر یقینی بناندی اے کہ~$d$
  دا آخرلا !!{formula} واقعی
  جملہ اے۔ آخری سطر اوس صورت
  لئی اے جدوں $d$ بس اک
  مفروضہ ہووے۔
\end{proof}

\begin{prob}
  ہِٹھلیاں خاصیتاں
  \olref[inc][art][pnd]{prop:followsby}
  دی طرز اُتے تعریف کرو:
  \begin{enumerate}
  \item $\fn{FollowsBy}_{\Elim{\lif}}(d)$،
  \item $\fn{FollowsBy}_{\Elim{\eq}}(d)$،
  \item $\fn{FollowsBy}_{\Elim{\lor}}(d)$،
  \item $\fn{FollowsBy}_{\Intro{\lforall}}(d)$۔
  \end{enumerate}
  آخری صورت لئی ایہ وی وکھاؤ کہ گوڈل
  نمبر~$d$ والے !!{derivation}
  دا آخری استنباط ایگن متغیر
  والی شرط پوری کردا اے یا
  نہیں، ایہ بنیادی بازگشتی
  ڈھنگ نال پرکھیا جا سکدا اے۔
  یعنی ایگن متغیر~$a$،
  جیہڑا $\Intro{\lforall}$
  استنباط دا اے،
  نہ~$d$ دے آخرلے
  !!{formula} وچ آوے،
  نہ~$d$ دے کسے
  خارج نہ کیتے مفروضے
  وچ۔ ایس لئی تسیں
  \olref[inc][art][pnd]{prop:openassum}
  والا بنیادی بازگشتی
  محمول $\fn{OpenAssum}$
  ورت سکدے او۔
\end{prob}

\begin{prop}
  \ollabel{prop:deriv}
  تعلق $\fn{Deriv}(d)$،
  جیہڑا اوس ویلے سچ
  ہوندا اے جدوں $d$
  کسے درست
  !!{derivation}~$\delta$
  دا گوڈل نمبر ہووے،
  بنیادی بازگشتی اے۔
\end{prop}

\begin{proof}
  !!^a{derivation}~$\delta$
  اوس ویلے درست اے جدوں
  اوہدا ہر استنباط کسے
  قاعدے دا درست اطلاق
  ہووے؛ یعنی اوہدے ذیلی
  !!{derivation}s وچوں
  ہر اک درست استنباط
  اُتے ختم ہووے۔ ایس
  لئی $\fn{Deriv}(d)$
  عین اوس ویلے سچ
  اے جدوں
  \[
  \bforall{i<\len{\fn{SubtreeSeq}(d)}}{\fn{Correct}((\fn{SubtreeSeq}(d))_i)}
  \]
\end{proof}

\begin{prop}
  \ollabel{prop:openassum}
  تعلق $\fn{OpenAssum}(z, d)$ بنیادی بازگشتی اے۔ ایہ قائم اے جے
$z$ کسے !!a{undischarged} مفروضے~$!A$ دا گوڈل نمبر ہووے جیہڑا
!!{derivation}~$\delta$ وچ اے، تے اوس اخذ دا گوڈل نمبر~$d$ ہووے۔
\end{prop}

\begin{proof}
  کسے مفروضے دی اک
  آمد اوس ویلے
  !!{discharged} ہوندی
  اے جدوں اوہ
  لیبل~$n$ نال
  $\delta$ دے کسے
  ذیلی !!{derivation}
  وچ آوے تے اوہ
  ذیلی اخذ اوسے لیبل~$n$ نال
  مفروضہ خارج کرن
  والے قاعدے اُتے
  ختم ہووے۔ سو
  $!A$، $\delta$
  دا !!a{undischarged}
  مفروضہ اوس ویلے
  اے جدوں اوہدیاں
  آمدواں وچوں گھٹ
  توں گھٹ اک
  $\delta$ وچ
  !!{discharged} نہ ہوئی
  ہووے۔ ایتھے ایہ گل یاد رکھو کہ
  $\delta$ وچ
  $!A$ دیاں
  !!{discharged} تے
  !!{undischarged}
  دونوں قسم دیاں
  آمدواں ہو سکدیاں
  نیں۔

  اک سلسلہ
  $\delta_0$، \dots،
  $\delta_k$ لو،
  جتھے $\delta_0 =
  \delta$، $\delta_k$
  مفروضہ
  $\Discharge{!A}{n}$
  ہووے (کسے~$n$ لئی)، تے
  $\delta_{i+1}$،
  $\delta_i$ دا
  براہِ راست ذیلی
  !!{derivation}
  ہووے۔ جے ایسا سلسلہ ملے جیہدے وچ کوئی
  وی~$\delta_i$
  لیبل~$n$ نال
  مفروضہ خارج
  کرن والے
  استنباط اُتے
  ختم نہ ہووے،
  تاں $!A$،
  $\delta$ دا
  !!a{undischarged}
  مفروضہ اے۔

  بنیادی بازگشتی فنکشن
  $\fn{SubtreeSeq}(d)$
  سانوں $\delta$ دے
  سارے ذیلی
  !!{derivation}s دے
  گوڈل نمبراں دا
  سلسلہ دیندا اے۔
  $\delta$ دے ذیلی
  !!{derivation}s دے
  گوڈل نمبراں دا
  ہر سلسلہ اوس دا
  ذیلی سلسلہ اے۔
  ذیلی سلسلہ ہون
  دا تعلق بنیادی
  بازگشتی اے:
  $\fn{Subseq}(s, s')$
  عین اوس ویلے
  سچ اے جدوں
  $\bforall{i<\len{s}}{\lexists[j<\len{s'}][(s)_i = (s')_j]}$۔
  براہِ راست ذیلی
  !!{derivation}
  ہون دا تعلق وی
  ایہوجیہا اے:
  $\fn{Subderiv}(d, d')$
  عین اوس ویلے
  سچ اے جدوں
  $\bexists{j<(d')_0}{d = (d')_j}$
  ہووے۔ سو $\fn{OpenAssum}(z, d)$ دی تعریف ایویں کر سکدے آں:
  \begin{multline*}
    \bexists{s<\fn{SubtreeSeq}(d)}{(\fn{Subseq}(s, \fn{SubtreeSeq}(d))
      \land (s)_0 = d \land {}} \\
    \bexists{n<d}{((s)_{\len{s} \tsub 1} = \tuple{0, z, n} \land {}}\\
      \bforall{i<(\len{s} \tsub 1)}{(\fn{Subderiv}((s)_{i+1}, (s)_i) \land {}}\\
      \fn{DischargeLabel}((s)_i) \neq n))).
  \end{multline*}
\end{proof}

\begin{prop}
  \ollabel{prop:prf-prim-rec}
  فرض کرو $\Gamma$،
  !!{sentence}s دا
  بنیادی بازگشتی
  سیٹ اے۔ تاں
  تعلق
  $\Prf[\Gamma](x, y)$،
  جیہڑا ایہ بیان
  کردا اے کہ
  «$x$، کسے
  !!a{derivation}~$\delta$
  دا کوڈ اے،
  جیہڑا $!A$ نوں $\Gamma$ وچلے !!{undischarged} مفروضیاں توں کڈھدا اے، تے
  $y$،~$!A$
  دا گوڈل نمبر
  اے»، بنیادی
  بازگشتی اے۔
\end{prop}

\begin{proof}
  فرض کرو
  «$y \in \Gamma$»
  بنیادی بازگشتی محمول~$R_\Gamma(y)$ نال دِتّا گیا اے۔
  سانوں دکھانا اے
  کہ تعلق
  $\Prf[\Gamma](x, y)$
  بنیادی بازگشتی
  اے۔ ایہ عین
  اوس ویلے سچ
  ہوندا اے جدوں
  $y$ کسے
  جملے~$!A$ دا
  گوڈل نمبر ہووے،
  تے $x$ کسے
  فطری استنباطی
  !!{derivation}
  دا کوڈ ہووے
  جیہدا آخرلا
  !!{formula}~$!A$
  اے، تے جیہدے
  سارے !!{undischarged}
  مفروضے
  $\Gamma$ وچ
  ہون۔

  \olref{prop:deriv}
  دے مطابق
  خاصیت
  $\fn{Deriv}(x)$،
  جیہڑی عین اوس
  ویلے سچ اے
  جدوں $x$
  فطری استنباط
  وچ کسے درست
  !!{derivation}~$\delta$
  دا گوڈل نمبر
  ہووے، بنیادی
  بازگشتی اے۔
  ایس لئی اسیں
  $\Prf[\Gamma](x, y)$ دی تعریف ایویں کر سکدے آں:
  \begin{align*}
    \Prf[\Gamma](x, y) \defiff {}
    & \fn{Deriv}(x) \land \fn{EndFmla}(x) = y \land {} \\
    & \bforall{z < x}{(\fn{OpenAssum}(z, x) \lif R_\Gamma(z))}.
  \end{align*}
\end{proof}

\end{document}
